% S04-M1 (package abhijit/s04-math): standing notation, the ordinary projector, Lemma ekills.
% Labels: ss:s04-notation, def:s04-e, lem:s04-ekills, if:s04-ekills.
\section{The ordinary projector and Lemma ekills}

\subsection{Standing notation}\label{ss:s04-notation}
Throughout stage~04, $p\ge5$ and $p\nmid D$. We write $\cO=\cOE$ for the ring of integers of a finite extension $E/\Qp$, with
uniformiser $\varpi$, and $\Lam=\cOE[[1+p\Zp]]$; for $u\in1+p\Zp$, $[u]\in\Lam$ is the group element. A \emph{level prime to $p$}
is a compact open subgroup $U=U^{(p)}\times\GL_2(\Zp)\subseteq\hat\cO^\times$ (read at $p$ through $\iota_p$) with $\Nrd U=\Zhat^\times$.
For such $U$ and $s\ge0$ put
\[UU_s:=U^{(p)}\times\Kp{s},\qquad \Kp{s}:=\iota_p^{-1}\Uone{p^s}_p,\qquad \Kpp{s}:=\Kp{s}\cap\iota_p^{-1}\Uzero{p^{s+1}}_p,\]
so that $UU_s=\Uone{Mp^s}$ if $U=\Uone{M}$ (in the notation of stage~01, $X_s(\Uone M)$ is the class set written $S_s$ there). Put
\[X_s(U):=\XB{UU_s},\qquad X'_s(U):=\XB{U^{(p)}\Kpp{s}},\]
and write $X_s$, $X'_s$ when $U$ is fixed. $\pi=\pi_{s+1,s}:X_{s+1}\to X_s$ is the natural projection, $\pi_*$ the pushforward on
$\Zp[X_{s+1}]\to\Zp[X_s]$ and $\pi^*$ the pullback of functions.

\emph{Facts used throughout} (all from stage~01). Let $s\ge1$.
\begin{enumerate}[label=(N\arabic*)]
\item Every $\Gamma_x$ at the levels $UU_s$, $U^{(p)}\Kpp s$ and $\Uzero qUU_s$ ($q\ne p$) is trivial, since these levels are contained in
  $\Uone p$ at $p$ (Lemma~\ref{lem:s01-stab}(b)). In particular $-1\notin UU_s$, $w_x=1$, the Brandt pairing is
  $\langle[x],[y]\rangle=\delta_{xy}$, and the norm map $\nu:\Zp[X_s]\to\Zp^{X_s}$ is $[x]\mapsto\mathbf 1_x$
  (Def.~\ref{def:s01-pairing}).
\item $\Delta_s:=(1+p\Zp)/(1+p^s\Zp)$ acts on $X_s$ by adelic diamonds $\langle d\rangle[x]=[x\delta_d]$, $\delta_d=\iota_p^{-1}\mathrm{diag}(1,d)$,
  and this action is free (Lemma~\ref{lem:s01-stab}(b)). The full group $(\Z/p^s)^\times=\mu_{p-1}\times\Delta_s$ acts as well, but not
  necessarily freely; only $\Delta_s$ is used for freeness statements.
\item \emph{Fibres.} If $U''\subseteq U'$ are levels with trivial stabilisers at $U'$, the fibre of $\XB{U''}\to\XB{U'}$ over $[x]$ is
  $\Gamma_x\backslash U'/U''=U'/U''$, so every fibre has exactly $[U':U'']$ elements. With $[\Kp s:\Kp{s+1}]=p^2$
  (Lemma~\ref{lem:s01-Up}(c)) this gives $|X_{s+1}|=p^2|X_s|$, hence $|X_s|=p^{2(s-1)}|X_1|$ and $v_p|X_s|\ge2s-2$.
\item $U_p$ is the operator $\NUp$ of Def.~\ref{def:s01-NUp}: $U_p[x]=\sum_{j=0}^{p-1}[x\varpi_j]$ on $\Zp[X]$, and
  $(U_pf)(x)=\sum_j f(x\varpi_j)$ on $\Zp^X$ (\emph{cohomological} $U_p$), with
  $\varpi_j=\iota_p^{-1}\left(\begin{smallmatrix}p&j\\0&1\end{smallmatrix}\right)$. The homological operator is the transpose of the
  cohomological one for the naive evaluation pairing $\Zp[X]\times\Zp^X\to\Zp$, $([x],f)\mapsto f(x)$. It commutes with $\pi_*$
  (Lemma~\ref{lem:s01-Up}(b)) and with $\langle d\rangle$, $T_\ell$, $S_\ell$, $\Pi_\ell$ and right translations by $B_q^\times$, $q\ne p$
  (Lemma~\ref{lem:s01-Up}(e), Prop.~\ref{prop:s01-relations}(c)).
\end{enumerate}
All modules are homological ($\Zp[X]$ with pushforward) unless marked otherwise. Interfaces of stages 02, 03, 05 and 06 are cited
by their names \texttt{if:s0N-\dots}; they are not proved in this chapter.

\subsection{The ordinary projector}

\begin{definition}[ordinary projector; seed \S2.6]\label{def:s04-e}
For a finitely generated $\Zp$-module $M$ with an endomorphism $T$, $e_T:=\lim_{n\to\infty}T^{n!}$. It exists, is an idempotent
in the closed subalgebra $\Zp[[T]]\subseteq\End M$ generated by $T$, and $M=e_TM\oplus(1-e_T)M$ with $T$ invertible on $e_TM$ and
topologically nilpotent on $(1-e_T)M$. The functor $M\mapsto e_TM$ is exact on $\Zp[T]$-modules of this kind. If $M$ is finite,
$e_T=T^{N}$ for any $N$ that is a multiple of the period and at least the index of $T$ in the finite monoid $\End M$
(e.g.\ $N=|\End M|!$). We write $e:=e_{U_p}$ and $e^*:=e_{U_p^*}$, where $U_p^*$ is the $\langle,\rangle$-adjoint of $U_p$; at level
$\Uone{Mp^s}$, $U_p^*=\langle p\rangle_MU'_p$ (Lemma~\ref{lem:s01-Up}(d), \pathtt{if:s01-hecke}), and Lemma~\ref{lem:s04-ekills}(d)
below gives the same formula at every level $UU_s$.
\end{definition}
\noindent\stproved{} below. The existence of the limit is also \stknown{}
\lit{hida\_1993\_elementary-theory-of-l-functions-and-eisenstein-series}{\S7.2, Lemma~1, p.~201: for a commutative $\cO$-algebra $A$
of finite rank and $x\in A$, $\lim x^{n!}$ exists and is idempotent} (apply it to $A=\Zp[T]\subseteq\End M$ when $M$ is free).
The proof below is elementary and also covers $M$ with torsion. \emph{Lean:} \stformal{} in the form of design D-e (Step~5), for $\cO=\Zp$ and unconditionally: \lean{Langlandscape.Ordinary.ordPart},
\lean{Langlandscape.Ordinary.nilPart}, \lean{Langlandscape.Ordinary.HasOrdDecomp} (their complementarity), the projector
\lean{Langlandscape.Ordinary.ordProj} with \lean{Langlandscape.Ordinary.ordProj\_idem}, \lean{Langlandscape.Ordinary.ker\_ordProj}, functoriality
\lean{Langlandscape.Ordinary.ordProj\_map}, exactness \lean{Langlandscape.Ordinary.ordPart\_exact}, and the existence for every f.g.\ $\Zp$-module
\lean{Langlandscape.Ordinary.padicOrdDecomp\_holds} (this discharges \lean{Langlandscape.Hyp.PadicOrdDecomp}; no hypothesis remains). Not formalized: the
limit formula $e_T=\lim T^{n!}$ itself and general $\cO_E$ (the finite case is covered by \lean{Langlandscape.Ordinary.hasOrdDecomp\_of\_finite}).

\begin{proof}
\emph{Step 1: finite monoids.} Let $x$ be an element of a finite monoid $S$. The powers $x,x^2,\dots$ take at most $|S|$ values, so
there are a least $i\ge1$ (the \emph{index}) and a least $\pi\ge1$ (the \emph{period}) with $x^{i+\pi}=x^i$; then $x^{a+\pi}=x^a$ for
all $a\ge i$, and $C:=\{x^a:a\ge i\}$ has exactly $\pi$ elements. For $a,b\ge i$, $x^ax^b=x^{a+b}$, and $a\mapsto x^a$ identifies $C$
with $\Z/\pi$; so $C$ is a cyclic group, and its identity is $x^N$ for every $N\ge i$ with $\pi\mid N$. Its identity is the only
idempotent in $C$. As $i,\pi\le|S|$, every $n\ge|S|$ has $n!\ge i$ and $\pi\mid n!$, so $x^{n!}=x^N$ for all such $n$; in
particular $N=|S|!$ works.

\emph{Step 2: finite $M$.} Apply Step~1 to $T$ in $S=\End M$ and put $e_T:=T^N$. Then $e_T$ is an idempotent and a polynomial in $T$,
so $M=e_TM\oplus(1-e_T)M$ with both summands $T$-stable. On $e_TM$, $T\cdot T^{N-1}=T^N=\mathrm{id}$ (take $N\ge2$), so $T$ is
invertible there. On $(1-e_T)M$, $T^N=T^N(1-T^N)=T^N-T^{2N}=0$ (as $T^{2N}=T^N$), so $T$ is nilpotent there. \emph{Uniqueness:} if
$M=M_1\oplus M_2$ with $T$-stable summands, $T$ invertible on $M_1$ and nilpotent on $M_2$, then $T^N|_{M_1}$ is an invertible
idempotent, hence the identity, and $T^N|_{M_2}$ is a nilpotent idempotent, hence $0$. So $e_T$ is the projection onto $M_1$
along $M_2$, i.e.\ $M_1=e_TM$ and $M_2=(1-e_T)M$. \emph{Functoriality:} if $f:M\to M'$ satisfies $fT=T'f$, then $fT^n=T'^nf$ for all
$n$, and with $N$ a common value for $(M,T)$ and $(M',T')$ we get $fe_T=e_{T'}f$.

\emph{Step 3: finitely generated $M$.} $M$ is $p$-adically complete and separated, $M=\varprojlim_mM/p^mM$, and each $M/p^mM$ is
finite. Let $e_m$ be the idempotent of Step~2 for $T$ on $M/p^mM$. The reductions $M/p^{m+1}M\to M/p^mM$ are $T$-equivariant, so
by functoriality the $e_m$ are compatible, and $e_T:=\varprojlim_me_m\in\End M$ is an idempotent commuting with $T$. For each $m$
there is $n_m$ with $T^{n!}\equiv e_T\pmod{p^m}$ on $M$ (i.e.\ $(T^{n!}-e_T)M\subseteq p^mM$) for all $n\ge n_m$; this is the
convergence $T^{n!}\to e_T$, and it shows that $e_T$ lies in the closure of $\Zp[T]$. $T$ is invertible on $e_TM$: it is injective,
because $Tx=0$ and $x=e_Tx$ give $x=\lim T^{n!}x=0$; and it is surjective, because $T(e_TM)+p\,e_TM=e_TM$ (since
$e_TM/p\,e_TM=e_1(M/pM)$ and $T$ is invertible there), so $T(e_TM)=e_TM$ by Nakayama. $T$ is topologically nilpotent on
$(1-e_T)M$: on $(1-e_T)M/p^m$ it is nilpotent by Step~2, so $T^n((1-e_T)M)\subseteq p^mM$ for $n\ge|\End(M/p^mM)|$.
\emph{Uniqueness} holds as in Step~2: if $M=M_1\oplus M_2$ ($T$-stable, $T$ invertible on $M_1$, topologically nilpotent on $M_2$),
reduce modulo $p^m$ and apply Step~2 to each $m$; the two decompositions agree modulo every $p^m$, hence agree. \emph{Functoriality}
as in Step~2: a $\Zp$-linear $f$ with $fT=T'f$ satisfies $f(p^mM)\subseteq p^mM'$, so $fe_T=e_{T'}f$ by passing to the limit.

\emph{Step 4: exactness.} Let $0\to M_1\xrightarrow{f}M_2\xrightarrow{g}M_3\to0$ be exact, with endomorphisms $T_i$ and $f$, $g$
equivariant. By functoriality $f$, $g$ commute with the projectors, so the sequence is the direct sum of the two complexes
$0\to e M_1\to eM_2\to eM_3\to0$ and $0\to(1-e)M_1\to(1-e)M_2\to(1-e)M_3\to0$. A direct summand of an exact complex is exact.

\emph{Step 5: the form used in Lean (D-e).} Put $\mathrm{ordPart}:=\bigcap_nT^nM$ and $\mathrm{nilPart}:=\{m:\forall k\ \exists n,\
T^nm\in p^kM\}$. Then $\mathrm{ordPart}=e_TM$ and $\mathrm{nilPart}=(1-e_T)M$. Indeed $e_TM\subseteq T^n(e_TM)$ for all $n$, and
$T^n((1-e_T)M)\subseteq p^kM$ for $n$ large, so $\bigcap_nT^nM\subseteq e_TM+\bigcap_kp^kM=e_TM$ (Krull: $\bigcap_kp^kM=0$). If $m$
lies in $\mathrm{nilPart}$, then $T^n(e_Tm)=e_T(T^nm)\in p^k\,e_TM$ for $n$ large, and since $T$ is invertible on $e_TM$ and
$T^{-1}$ preserves $p^ke_TM$, $e_Tm\in p^ke_TM$ for all $k$, i.e.\ $e_Tm=0$; the converse inclusion is Step~3. So the two
submodules are complementary, which is the statement $\mathrm{HasOrdDecomp}\,T$ of the stage plan, and $e_T$ is the projection
onto the first along the second.
\end{proof}

\subsection{Lemma ekills}

\begin{lemma}[$e$ kills degree, vertex and component terms; seed Lemma ekills]\label{lem:s04-ekills}
Let $s\ge1$, $V=\Zp$, $U$ of level prime to $p$.
\begin{enumerate}[label=(\alph*)]
\item $\deg:\Zp[X_s(U)]\to\Zp$ satisfies $\deg\circ U_p=p\cdot\deg$. Hence $e$ kills $\Zp$ (with $U_p=p$), $e\,\Zp[X_s]^0=e\,\Zp[X_s]$,
  and, for $q\nmid Dp$ and $U_q=\GL_2(\Zq)$, $e$ kills $H_0(\Gamma_{UU_s},\Zp)\cong\Zp$ (with the $U_p$-action induced from
  $\coker\partial$, \texttt{if:s02-coh}).
  \provisional{S04-M1: the $H_0$ clause uses that $\BG(UU_s,q)$ is connected (as in Thm.~lambdaribet, F7); the stage tex's proof
  says ``the connected graph'' but the statement does not list the hypothesis.} Mathematically the hypothesis is a theorem:
  $\Nrd(UU_s)=\hat\Z^\times$ and $B_q$ is split, so strong approximation for $B^1$ (Eichler, Kneser) gives connectivity; it is kept as
  a hypothesis because the Lean side carries it as \texttt{Hyp.SA}. (Also, $H_0(\Gamma,\Zp)\cong\Zp$ holds for any group; connectivity
  is needed for its identification with $\coker\partial$.)
\item ($B=B^{(q)}$, the definite algebra ramified exactly at $\{q,\infty\}$, written $B_q$ in stage~03, Def.~\ref{def:s03-DR}; not
  the local algebra $B\otimes\Q_q$ of (N4); here $U=\Uone M$, so that $UU_s=\Uone{Mp^s}$ is a level of Def.~\ref{def:s03-DR}, defined on
  $\hat B^{(q)\times}$ by the same local conditions; DR graph $\GDR_q(UU_s)$, \pathtt{if:s03-graph}.) Let $U_p$ act on the vertex module $\Zp[\{v_0,v_1\}]$ by $p$; then
  $\partial$ is $U_p$-equivariant, $e$ kills the vertex term, and $e\,\Xcomb_q(UU_s)\otimes\Zp=e\,\Zp[\XBq{UU_s}]$.
\item ($B=B^{(q)}$ as in (b).) $\Phicomb(\GDR_q(UU_s))\otimes\Zp\cong\Zp/|X_s|$ carries the action $\varphi\mapsto\varphi\circ U_p^*$ induced from
  $\Hom(X,\Z)$, which is multiplication by $p$; so $e(\Phicomb\otimes\Zp)=0$. This is a non-trivial statement: $v_p|X_s|\ge 2s-2$.
\item (Any $B$.) The Brandt pairing $\langle[x],[y]\rangle=\delta_{xy}$ on $\Zp[X_s]$ restricts to a perfect pairing
  \[e\,\Zp[X_s]\times e^*\Zp[X_s]\to\Zp,\]
  and $\nu$ identifies $e^*\Zp[X_s]$ with $e\,\Zp^{X_s}$ for the \emph{cohomological} $U_p$ ($f\mapsto\sum_jf(\cdot\,\varpi_j)$).
  Thus the perfect pairing is the evaluation pairing $e\,\Zp[X_s]\times e\,\Zp^{X_s}\to\Zp$. ($e$ and $e^*$ differ in general:
  Remark~\ref{rem:s04-eestar}.)
\end{enumerate}
Here $U_p^*=T_{\varpi_p^{-1}}=\langle p\rangle_MU'_p$, where $U'_p=[UU_s\,\eta_p\,UU_s]$, $\eta_p=\iota_p^{-1}\mathrm{diag}(1,p)$, and
$\langle p\rangle_M$ is right translation by the central idele $t$ with $t_\ell=p$ for $\ell\ne p$ and $t_p=1$; for
$U^{(p)}=\Uone M$ this is the diamond $\langle d\rangle$ with $d\equiv p\bmod M$, $d\equiv1\bmod p^s$ (Lemma~\ref{lem:s01-Up}(d)).
\stproved{} below ((a)--(c) as in the seed, with the missing equivariance details; (d) with the identification of $e^*$; the
$H_0$ clause of (a) under the connectivity hypothesis). \emph{Lean} (for $\Zp$, $U=\Uone M$, $p\ge5$; the argument \texttt{hdec} is discharged by
\pathtt{padicOrdDecomp\_holds}): (a), degree part, \stformal: \lean{Langlandscape.Ordinary.deg\_comp\_UpHom},
\lean{Langlandscape.Ordinary.ekillsA\_tower} ($e\,\Zp[X_s]^0=e\,\Zp[X_s]$; abstract form \lean{Langlandscape.Ordinary.ordPart\_ker\_deg}); (d), perfect pairing,
\stformal: \lean{Langlandscape.Ordinary.perfect\_ordPart\_pairing\_tower\_padic}, \lean{Langlandscape.Ordinary.brandtPairing\_eq\_deltaPairing}. For (b), (c) only
the \emph{abstract} parts are formalized (\lean{Langlandscape.Ordinary.ordPart\_le\_ker\_of\_vertex}, \lean{Langlandscape.Ordinary.ordPart\_tensor\_phiComb\_eq\_bot}, with the
graph inputs and the identification $\Phicomb\cong\Z/|X_s|$ as hypotheses; the valuation bound $v_p|X_s|\ge2s-2$ modulo \pathtt{Hyp.EichlerMass} and a level-index input
\texttt{hlev}). Not formalized: the $H_0$ clause, the flatness identification $\Xcomb\otimes\Zp=\ker\partial_{\Zp}$, and the identification of $e^*$ with $e\,\Zp^{X_s}$
under $\nu$. So the lemma as a whole stays \stproved{} (tex).
\end{lemma}
\begin{proof}
(a) $\deg U_p[x]=\#\{j\}=p$, so $\deg\circ U_p=p\deg$. On $\Zp$ with $U_p=p$, $p^{n!}\to0$, so $e=0$. Since $\deg$ is $U_p$-equivariant
for this action, Def.~\ref{def:s04-e} (exactness) applied to $0\to\Zp[X_s]^0\to\Zp[X_s]\xrightarrow{\deg}\Zp\to0$ gives
$0\to e\Zp[X_s]^0\to e\Zp[X_s]\to0$, i.e.\ $e\Zp[X_s]^0=e\Zp[X_s]$.
For $H_0$: the level $UU_s$ is free (N1), and when $\BG(UU_s,q)$ is connected (equivalently Prop.~quot holds at this level,
\texttt{if:s02-quot}), $H_0(\Gamma_{UU_s},\Zp)$ is the cokernel of the boundary map $\partial$ of the two-term complex of
\texttt{if:s02-coh}, whose target is $\Zp[X_s]^2$ (the two vertex orbits) and whose source is $\Zp[\XB{\Uzero qUU_s}]$ (the edges),
$\partial=(-s_{0*},s_{1*})$. The map $(a,b)\mapsto\deg a+\deg b$ kills $\partial$ of every edge ($-1+1=0$) and induces
$\coker\partial\cong H_0(|\BG|)\otimes\Zp=\Zp$ because the graph is connected. $U_p$ acts componentwise on $\Zp[X_s]^2$; it commutes
with $s_{0*},s_{1*}$ (\texttt{if:s02-quot}(c); both are supported at $q\ne p$, cf.\ (N4)), so it acts on $\coker\partial$, and there by
$p$, since $\deg\circ U_p=p\deg$ on each component. Hence $e$ kills $H_0$. No density statement is used: only connectivity of the
quotient graph.

(b) In $\GDR_q(UU_s)$ every edge $\epsilon\in E=\XBq{UU_s}$ runs from $v_0$ to $v_1$ and all weights are $1$ (N1), so
$\partial\epsilon=v_1-v_0$ and $\partial\big(\sum a_\epsilon\epsilon\big)=\deg(a)\,(v_1-v_0)$. Hence
$\partial(U_p\epsilon)=\sum_j\partial[\epsilon\varpi_j]=p(v_1-v_0)=p\,\partial\epsilon$, so $\partial$ is $U_p$-equivariant when $U_p$
acts on $\Z[\{v_0,v_1\}]$ by $p$; on the vertex module $e=0$ as in (a). The sequence $0\to\Xcomb_q\to\Z[E]\xrightarrow{\partial}
\Z[V]^0\to0$ is exact with $\Xcomb_q=\Z[E]^0$ (\pathtt{if:s03-graph}), and $\Z[V]^0=\Z(v_1-v_0)\cong\Z$ with $U_p=p$. Tensor with
$\Zp$ and apply $e$: $e\,\Xcomb_q(UU_s)\otimes\Zp\xrightarrow{\sim}e\Zp[E]$.

(c) With $w=1$ the lattice is $X=\Z[E]^0$, $\Hom(X,\Z)=\Z^E/\Z\mathbf 1$ (restriction from $\Z^E$; the kernel consists of the
constants), and $X\to\Hom(X,\Z)$, $x\mapsto\langle x,\cdot\rangle$, is the inclusion $\Z[E]^0\subseteq\Z[E]=\Z^E$ followed by the
projection. So $\Phicomb=\Z^E/(\Z[E]^0+\Z\mathbf 1)$, and $\Sigma:\varphi\mapsto\sum_\epsilon\varphi(\epsilon)$ induces
$\Phicomb\cong\Z/|E|$ (as $\Sigma$ identifies $\Z^E/\Z[E]^0$ with $\Z$ and sends $\mathbf1$ to $|E|$). The Hecke action:
$\langle U_px,y\rangle=\langle x,U_p^*y\rangle$, so $x\mapsto\langle x,\cdot\rangle$ intertwines $U_p$ on $X$ with
$\varphi\mapsto\varphi\circ U_p^*$ on $\Hom(X,\Z)$; both $U_p$ and $U_p^*$ preserve $\Z[E]^0$, since $U_p^*=T_{\varpi_p^{-1}}$ also has
$p$ cosets (its double coset is $z_p^{-1}U\eta_pU$, Lemma~\ref{lem:s01-Up}(c),(d)). On $\Z/|E|$:
\[\Sigma(\varphi\circ U_p^*)=\sum_\epsilon\varphi(U_p^*\epsilon)=\sum_{\epsilon'}\varphi(\epsilon')\cdot c_{\epsilon'},\qquad
c_{\epsilon'}=\sum_\epsilon(\text{coefficient of }\epsilon'\text{ in }U_p^*\epsilon).\]
With $w=1$ the matrix of $U_p^*$ is the transpose of that of $U_p$, so $c_{\epsilon'}$ is the column sum of $U_p$ at $\epsilon'$, i.e.\
$\deg U_p[\epsilon']=p$. So $U_p$ acts on $\Phicomb\cong\Z/|E|$ by $p$, and on $\Phicomb\otimes\Zp=\Zp/p^{v_p|E|}$ by the nilpotent
map $p$; hence $e=0$ there. Non-triviality: $|E|=|\XBq{UU_s}|=p^{2(s-1)}|\XBq{UU_1}|$ by (N3), so $v_p|E|\ge2s-2$. (For $U=\Uone M$
this agrees with the mass formula of \texttt{if:s01-mass}: $\psi(Mp^s)\varphi(Mp^s)=p^{2(s-1)}\psi(Mp)\varphi(Mp)$, Thm.~\ref{thm:s01-mass}.)

(d) \emph{The adjoint.} By (N1) the pairing is $\delta_{xy}$, symmetric and perfect on the finite free module $\Zp[X_s]$. By Step~1 of
the proof of Lemma~\ref{lem:s01-Up}(d) (orbit counting, valid for any level and trivial coefficients), the adjoint of
$T_g=[UU_s\,g\,UU_s]$ is $T_{g^{-1}}$; so $U_p^*=T_{\varpi_p^{-1}}$. Steps~2--3 of that proof use only that $p\in\Q^\times$ is central
and that $t$ is central, hence normalises $UU_s$; they give $U_p^*[y]=\sum_j[y\eta^{(s)}_jt]=\langle p\rangle_MU'_p[y]$ at every level
$UU_s$, with $\langle p\rangle_M[y]:=[yt]$. \emph{Perfectness.} From $\langle U_p^{n!}a,b\rangle=\langle a,U_p^{*\,n!}b\rangle$ and
continuity of the pairing, $\langle ea,b\rangle=\langle a,e^*b\rangle$. Hence $\langle ea,(1-e^*)b\rangle=\langle a,e^*(1-e^*)b\rangle=0$ and
$\langle(1-e)a,e^*b\rangle=\langle e(1-e)a,b\rangle=0$. So with respect to the decompositions $\Zp[X]=e\Zp[X]\oplus(1-e)\Zp[X]$ (first
variable) and $\Zp[X]=e^*\Zp[X]\oplus(1-e^*)\Zp[X]$ (second variable) the pairing is block diagonal, and the isomorphism
$\Zp[X]\to\Hom(\Zp[X],\Zp)$, $a\mapsto\langle a,\cdot\rangle$, is the direct sum of $e\Zp[X]\to\Hom(e^*\Zp[X],\Zp)$ and
$(1-e)\Zp[X]\to\Hom((1-e^*)\Zp[X],\Zp)$. A direct sum of two maps is bijective iff both are, so the first is an isomorphism:
the pairing $e\Zp[X]\times e^*\Zp[X]\to\Zp$ is perfect. \emph{Identification.} $\nu(a)(y)=\langle a,[y]\rangle$, and for the
cohomological $U_p$,
\[(U_p\,\nu(a))(y)=\sum_j\langle a,[y\varpi_j]\rangle=\langle a,U_p[y]\rangle=\langle U_p^*a,[y]\rangle=\nu(U_p^*a)(y),\]
so $\nu\circ U_p^*=U_p^{\rm coh}\circ\nu$ and, by functoriality (Def.~\ref{def:s04-e}), $\nu e^*=e^{\rm coh}\nu$. Thus $\nu$ maps
$e^*\Zp[X_s]$ onto $e\,\Zp^{X_s}$, and $\langle a,b\rangle=\nu(b)(a)$ turns the perfect pairing into the evaluation pairing
$e\Zp[X_s]\times e\Zp^{X_s}\to\Zp$.
\end{proof}
\begin{remark}[$e\ne e^*$]\label{rem:s04-eestar}
By (d), $e^*$ is the adjoint of $e$ for $\langle[x],[y]\rangle=\delta_{xy}$, so $e=e^*$ iff $e$ is self-adjoint. This fails in general:
for $D=2$ (Hurwitz order), $M=7$, $p=5$, $s=1$ ($48$ points, all $\Gamma_x$ trivial, so the adjoint is the transpose) the idempotent
$e=e_{U_p}$ has rank $27$ and $e\not\equiv e^{\mathsf T}=e^*\pmod 5$. \stcomp{} (script
\pathtt{code/python/langlandscape/ordinary/e\_vs\_estar.py}, output \pathtt{code/results/s04/M95\_e\_vs\_estar.json}; the matrix of
$U_p$ is that of the S04-C3a sign check, \pathtt{code/results/s04/H6\_sign.json}.)
\end{remark}

\noindent\emph{Why (d) matters.} The family $(e\Zp[X_s])_s$ is an inverse system under $\pi_*$ (Lemma~\ref{lem:s01-Up}(b)), and
$(e\Zp^{X_s})_s$ is a direct system under $\pi^*$. Via $\nu$, (d) identifies $e^*\Zp[X_s]$ with $e\Zp^{X_s}$, and on $\Zp[X_s]$ the
direct system $\pi^*$ becomes the transfer $\pi^!$ ($\nu'\circ\pi^!=\pi^*\circ\nu$ inside the tower, Prop.~\ref{prop:s01-adjoint}(c)); so
$(e^*\Zp[X_s])_s$ is a direct system under $\pi^!$. It is not stable under $\pi_*$ ($U'_p$ does not commute with $\pi_*$). Thus (d)
pairs an inverse system $(e,\pi_*)$ with a direct system $(e^*,\pi^!)$, and the $\Lambda$-adic pairing cannot be the naive limit of (d);
it needs the twist of Def.~\ref{def:s04-BN}.

\noindent\emph{Right translations.} Lemma~\ref{lem:s04-ekills} uses right translations only through the commutation of
$U_p$ with $s_{1*}$ (right translation by an element of $B_q^\times$). We use the direction fixed in
Prop.~\ref{prop:s01-relations}(c): $R_gf=f(\cdot\,g)$ maps $\Brandt{g^{-1}Ug}{V}$ to $\Brandt{U}{V}$, and on class sets
$\rho_g[x]=[xg]$ maps $\XB U\to\XB{g^{-1}Ug}$; both commute with $U_p$ since $g$ is supported at $q\ne p$.

\begin{interface}[\pathtt{if:s04-ekills}]\label{if:s04-ekills}
Lemma~\ref{lem:s04-ekills} (a)--(c), plus (d): perfect pairing $e\Zp[X_s]\times e^*\Zp[X_s]\to\Zp$, equivalently the evaluation pairing
$e\Zp[X_s]\times e\Zp^{X_s}$ (cohomological $e$). Status \stproved{} (the $H_0$ clause of (a) under connectivity of $\BG(UU_s,q)$).
Addendum: Prop.~\ref{prop:s04-ekillsVk} ($e$ kills $H_0$ with $\Vk$), \sttoprove.
\end{interface}
